\chapter{DNA: the Molecule of Heredity}\label{ch:g9:dna}

\omterm{def:g9:dna:gene}{Genes} are made of \omterm{def:g9:dna:dna}{DNA}. Understanding \omterm{def:g9:dna:dna}{DNA}'s structure explains how
information is stored, copied and used --- the \omterm{def:g7:digestion:heart}{heart} of molecular
\omterm{def:g7:heredity:heredity}{heredity} at introductory level.

\section{DNA structure}

\begin{definition}[DNA]\label{def:g9:dna:dna}
\emph{DNA}\index{DNA} (deoxyribonucleic acid) is the molecule that stores
genetic information in \omterm{def:g7:microscope:cell}{cells}. It is a double helix of two strands.
\end{definition}

\omimg{dna-helix.jpg}{0.52\linewidth}{\omterm{def:g9:dna:dna}{DNA} double helix: two strands
twisted around each other.}

\begin{definition}[Nucleotides and bases]\label{def:g9:dna:base}
\omterm{def:g9:dna:dna}{DNA} is a polymer of \emph{nucleotides}\index{nucleotide}. Each nucleotide
has a sugar, a phosphate and a base. The four bases are adenine (A),
thymine (T), cytosine (C) and guanine (G).
\end{definition}

\begin{proposition}[Base pairing]\label{prop:g9:dna:pair}
Bases pair specifically: A with T, and C with G. The sequence of bases
along a strand encodes information; the complementary strand is fixed by
pairing rules.
\end{proposition}
\admitted

\begin{omfigure}
\begin{tikzpicture}[font=\footnotesize]
  \foreach \y/\L/\R in {0/A/T, 0.7/C/G, 1.4/T/A, 2.1/G/C, 2.8/A/T} {
    \node[draw, circle, fill=omDef!20, minimum size=0.55cm] at (0,\y) {\L};
    \node[draw, circle, fill=omMeth!20, minimum size=0.55cm] at (2.2,\y) {\R};
    \draw[thick, omProp] (0.3,\y) -- (1.9,\y);
  }
  \draw[very thick, omThm] (-0.5,-0.3) -- (-0.5,3.1);
  \draw[very thick, omThm] (2.7,-0.3) -- (2.7,3.1);
  \node at (1.1,-0.8) {complementary base pairs};
\end{tikzpicture}

{\small Ladder model: sides are sugar--phosphate backbones; rungs are base
pairs.}
\end{omfigure}

\section{Replication}

\begin{definition}[DNA replication]\label{def:g9:dna:rep}
\emph{Replication}\index{DNA replication} is the copying of \omterm{def:g9:dna:dna}{DNA} before
\omterm{def:g7:microscope:cell}{cell} division. Strands separate; each old strand templates a new
complementary strand --- \emph{semi-conservative}\index{semi-conservative
replication} copying.
\end{definition}

\begin{method}[Replication in three ideas]\label{met:g9:dna:rep}
\begin{enumerate}
  \item Unzip the double helix (break base pairs).
  \item Match free \omterm{def:g9:dna:base}{nucleotides} to each exposed strand (A--T, C--G).
  \item Two double helices result, each with one old and one new strand.
\end{enumerate}
\end{method}

\begin{omfigure}
\begin{tikzpicture}[font=\footnotesize]
  \draw[very thick, omDef] (0,0) -- (0,2.5);
  \draw[very thick, omMeth] (0.6,0) -- (0.6,2.5);
  \draw[-Stealth, thick] (1.2,1.2) -- (2.2,1.2);
  \draw[very thick, omDef] (3,0) -- (3,2.5);
  \draw[very thick, omExo] (3.6,0) -- (3.6,2.5);
  \draw[very thick, omMeth] (5.2,0) -- (5.2,2.5);
  \draw[very thick, omExo] (5.8,0) -- (5.8,2.5);
  \node at (0.3,-0.4) {parent};
  \node at (4.4,-0.4) {two copies};
\end{tikzpicture}

{\small \omterm{def:g9:dna:rep}{Semi-conservative} idea: each daughter \omterm{def:g9:dna:dna}{DNA} keeps one parental
strand.}
\end{omfigure}

\section{From gene to protein (preview)}

\begin{definition}[Gene as code]\label{def:g9:dna:gene}
A \emph{gene}\index{gene} is a \omterm{def:g9:dna:dna}{DNA} sequence that is used to build a
functional product, often a \omterm{def:g6:nutrition:groups}{protein}. The base sequence is a code for the
order of amino acids.
\end{definition}

\begin{proposition}[Central idea]\label{prop:g9:dna:central}
\omterm{def:g9:dna:dna}{DNA} is transcribed to an RNA message; \omterm{def:g8:organelles:ribosome}{ribosomes} translate that message
into \omterm{def:g6:nutrition:groups}{protein}. Details expand in high school; the key is: sequence
information flows from \omterm{def:g9:dna:dna}{DNA} to \omterm{def:g6:nutrition:groups}{proteins} that build and run the \omterm{def:g7:microscope:cell}{cell}.
\end{proposition}
\admitted

\begin{example}[Mutation]\label{ex:g9:dna:mut}
A \emph{mutation}\index{mutation} is a change in \omterm{def:g9:dna:dna}{DNA} sequence. Some change
\omterm{def:g6:nutrition:groups}{proteins} and \omterm{def:g7:heredity:trait}{traits}; many have little effect; some are harmful; rarely
helpful. \omterm{ex:g9:dna:mut}{Mutations} supply new \omterm{def:g5:evolution:variation}{variation} for \omterm{def:g8:evolution:def}{evolution}.
\end{example}

\section{Exercises}

\begin{exercise}[$\star$]\label{exo:g9:dna:1}
What does \omterm{def:g9:dna:dna}{DNA} stand for, and what is its main role?
\end{exercise}

\begin{exercise}[$\star$]\label{exo:g9:dna:2}
Name the four bases. Which pairs with which?
\end{exercise}

\begin{exercise}[$\star$]\label{exo:g9:dna:3}
If one strand is ATGCCA, what is the complementary strand?
\end{exercise}

\begin{exercise}[$\star$]\label{exo:g9:dna:4}
Why must \omterm{def:g9:dna:dna}{DNA} replicate before \omterm{def:g9:cell-division:mitosis}{mitosis}?
\end{exercise}

\begin{exercise}[$\star$]\label{exo:g9:dna:5}
What is a \omterm{def:g9:dna:gene}{gene} in terms of \omterm{def:g9:dna:dna}{DNA}?
\end{exercise}

